% ICAIF '26 Workshop Proposal — max 5 pages, 12pt, single PDF
% Submit via CMT3 "Workshop Proposals" track by July 19, 2026 23:59 AoE
\documentclass[12pt]{article}
\usepackage[margin=1in]{geometry}
\usepackage[hidelinks]{hyperref}
\usepackage{enumitem}
\setlist{nosep}

\title{\textbf{Workshop Proposal: Entity-Centric AI in Finance\\
\large Identity, Linkage, and Lifecycle Dynamics of Financial Entities}}
\author{}
\date{}

\begin{document}
\maketitle
\vspace{-3em}

\section{Workshop Topic and Title}

\textbf{Title:} \emph{The 1st Workshop on Entity-Centric AI in Finance
(WECAIF)} \hfill \textbf{Duration:} half day

Nearly every AI system in financial services rests on a layer that
receives far less research attention than the models above it: knowing
\emph{which real-world entity} a record refers to. Merchant and
counterparty resolution in payments, KYC/AML screening, vendor and
invoice deduplication, portfolio-company tracking, and financial
knowledge-graph construction all depend on entity identity --- and all
break in characteristic ways when entities merge, rebrand, spin off,
or migrate jurisdictions. This workshop convenes the researchers and
practitioners working on that layer: entity resolution and record
linkage for financial data, corporate-hierarchy and ownership graphs,
temporal/lifecycle dynamics of entities, open identifier
infrastructure (LEI, CIK, ISIN, OpenCorporates), and the emerging role
of LLMs as entity matchers and enrichers. The scope is squarely AI in
financial services: methods, benchmarks, datasets, and deployed-system
reports.

\section{Motivation}

Three developments make this topic timely for the ICAIF community.

\textbf{(i) Entity work has quietly become an LLM evaluation
frontier.} Recent large-scale studies of LLM-based entity matching
report near-saturated performance on classical point-in-time
benchmarks, while newly released lifecycle-aware benchmarks built on
open regulatory data (GLEIF, SEC EDGAR) show frontier models trailing
small fine-tuned encoders under uniform protocols. Reconciling these
findings --- what entity tasks are actually solved, which remain open,
and how evaluation design shapes the answer --- is exactly the kind of
methods-and-measurement conversation a focused workshop enables.

\textbf{(ii) Open identity infrastructure has reached critical mass.}
GLEIF's LEI corpus (3.3M+ entities, CC0, with structured lifecycle
events and relationship records), EDGAR, and Wikidata now support
fully reproducible research on problems that were previously locked
inside proprietary vendor datasets. A workshop can accelerate the
formation of shared benchmarks and data standards around this
infrastructure.

\textbf{(iii) Practitioner demand.} Entity mastering is a named
production concern at payment processors, spend-management platforms,
market-data vendors, and compliance teams, yet the practitioner and
academic communities rarely share a venue. ICAIF is the natural home
for that exchange; no prior ICAIF workshop has focused on the entity
layer (\S6).

The workshop aims to produce: a clearer map of open problems in
financial-entity intelligence, seeds for shared benchmarks and
evaluation protocols (including lifecycle-aware and
retrieval-grounded evaluation), and new academic--industry
collaborations.

\section{Organizing Committee}

% [TO FILL — names, affiliations, emails, 3–5 line bios each.
%  Reviewers weigh organizational experience heavily: include at least
%  two senior members with prior workshop organization if possible.]

\begin{itemize}
  \item \textbf{[Faculty co-chair — TO FILL]}, [affiliation, email].
        [Bio highlighting relevant research and any prior
        workshop/conference organization.] \emph{(primary contact)}
  \item \textbf{Sheng Chang Li}, Western University,
        \texttt{sli3244@uwo.ca}. Author of \textsc{Sylvester}, an open
        benchmark and method for lifecycle-aware entity resolution on
        GLEIF/EDGAR/Wikidata data (under review, ICAIF '26); builds
        entity-linkage tooling for brokerage-account data at
        SnapTrade. [Adjust/extend.]
  \item \textbf{[Senior co-organizer 2 — TO FILL]} --- ideally an
        industry researcher owning an entity-mastering system.
  \item \textbf{[Senior co-organizer 3 — TO FILL]} --- ideally a
        data-management/ER academic for cross-community reach.
\end{itemize}

\section{Tentative Program Committee and Invited Speakers}

% [TO FILL — the selection committee explicitly scores "quality of
%  proposed invited speakers". List only people you genuinely intend
%  to invite; mark confirmed vs. contacted vs. candidate.]

\textbf{Keynote candidates (to be confirmed):}
\begin{itemize}
  \item [Candidate A — e.g., a GLEIF or open-identifier-infrastructure
        leader, for the data/infrastructure perspective]
  \item [Candidate B — e.g., an author of large-scale LLM
        entity-matching studies, for the methods perspective]
\end{itemize}

\textbf{Invited talk candidates:} [3--4 names spanning: a production
entity-mastering system (payments or market data), financial
knowledge graphs, temporal entity linking, compliance/AML screening.]

\textbf{Tentative program committee:} [10--15 names across ER/data
management, NLP/IE, and quantitative finance; mix of academia and
industry. TO FILL.]

\section{Tentative Schedule and Format}

Half-day (four hours), hybrid-friendly but presented in person:

\begin{itemize}
  \item 00:00--00:10 \quad Opening: the entity layer in financial AI
  \item 00:10--00:45 \quad Keynote 1 (infrastructure/data)
  \item 00:45--01:45 \quad Contributed papers I (4 talks, 15 min each)
  \item 01:45--02:15 \quad Coffee + poster session
  \item 02:15--02:50 \quad Keynote 2 (methods/LLMs)
  \item 02:50--03:35 \quad Contributed papers II (3 talks, 15 min each)
  \item 03:35--04:00 \quad Panel: ``What would it take for a shared
        financial-entity benchmark suite?'' (keynotes + organizers +
        one PC member)
\end{itemize}

Contributions: short (4pp) and full (8pp) non-archival papers via a
lightweight CMT/OpenReview instance run by the organizers; posters for
work-in-progress and industry abstracts. The organizers will run the
call for papers, review (2 reviews per submission), and acceptance
decisions, and will maintain the workshop website.

\section{Comparable Workshops}

Finance-focused NLP and ML workshops --- e.g., the long-running
FinNLP series and AAAI's Knowledge Discovery from Financial Data
lineage, as well as ICAIF's own prior workshops on LLMs and
generative AI in finance --- have repeatedly featured entity-adjacent
papers (NER, KG construction, document intelligence), demonstrating
community supply. However, none centers the entity-identity layer
itself: resolution, linkage, lifecycle dynamics, and identifier
infrastructure. In the data-management community, ER workshops exist
but lack the financial-domain focus and practitioner participation
ICAIF provides. WECAIF is proposed as a new series (not a
continuation) occupying that intersection; its differentiation is the
exclusive focus on entity identity in financial services and its
benchmark/data-standards agenda.

\end{document}
